\scanpage{54}
par $E_{\bar K/K}$) en prenant son centralisateur $Z(L_i)$ dans $E$, d'où
\[
(2)\qquad
\diagram{\begin{tikzcd}[ampersand replacement=\&,column sep=small]
1\arrow[r,no head]\&L_i\arrow[r,no head]\arrow[d,phantom,"\simeq" description]\&Z(L_i)\arrow[r]\&\Gamma\arrow[r]\&1\\
\&T_\infty(\bar K^*)=T\&\&\&
\end{tikzcd}}
\]
et on prend les scindages de cette extension. Il en existe, p.ex. définis par une uniformisante de $\unsure{\mathcal O}_{\widehat U,i}$. L'ens des classes de conjugaison des \deleted{\unsure{ses}} scindages \deleted{\unsure{les}} de (\unsure{1|2}) est un torseur sous \deleted{$H^1(\Gamma,T)$}
\[
(3)\qquad H^1(\Gamma,T)\simeq\varprojlim_n H^1(K,{}_n\mu)\simeq\widehat{K^*}=\varprojlim_n K^*_n
\]
et \deleted{\unsure{cet ensemble}} s'envoie injectivement \deleted{\unsure{vers}} dans l'ens des classes de $\pi$-conj. de scindages.

\uline{Proposition} On suppose $(g,\nu)\ne(0,0),(0,1)$ i.e. $\Pi_{\bar U}=\Pi_{U/K}\ne(1)$. Alors les classes de \inserted{$\pi$-conjugaison} scindages \deleted{\unsure{assoc.}} de (1) définis par les scindages de (2) sont distinctes de celles associées aux pts de $U(K)$. \deleted{De plus,} Si \inserted{(de plus)} $(g,\nu)\ne(0,2)$, \deleted{\unsure{dans le}} \inserted{i.e. si on est} dans le cas anabélien, alors les classes de $\pi$-conjugaison de scindages de (1), associés à des scindages de (2) pour deux indices $i\deleted{\unsure{a}}=i_1$ et $i=i_2$ distincts \inserted{sont distinctes.}
\margin{C'est démontré sauf pour le type $(0,3)$, qui mérite un examen approfondi.}
\margin{53}
